\section*{Capítulo V. Sistemas de factores}\label{kapitel-v.-faktorensysteme}

\subsection*{§ 22. El grupo de los cuerpos con centro dado}\label{die-gruppe-der-koerper-mit-gegebenem-zentrum}

Fijemos como base de nuestras investigaciones un cuerpo conmutativo \(\mathsf P\). El conjunto de todos los anillos de matrices \(\mathfrak K_r\) cuyos cuerpos de automorfismos \(\mathfrak K\) son cuerpos de rango finito sobre \(\mathsf P\), con el propio \(\mathsf P\) como centro, es un sistema cerrado respecto de la ``formación del producto directo''. El producto directo de dos de tales anillos \(\mathfrak K_r\) y \(\mathfrak L_s\) es simplemente el anillo \(\mathfrak K_r\cdot\mathfrak L_s\) que se obtiene de \(\mathfrak K_r\) extendiendo el cuerpo base a \(\mathfrak L_s\):
\[
\mathfrak K_r\times\mathfrak L_s=\mathfrak K_r\cdot\mathfrak L_s.
\]
Según la observación de la página 22, esta formación del producto directo es conmutativa:
\[
\mathfrak K_r\times\mathfrak L_s=\mathfrak L_s\times\mathfrak K_r.
\]
La afirmación de que el sistema de todos los \(\mathfrak K_r\) es cerrado multiplicativamente no significa, pues, sino que \(\mathfrak K_r\times\mathfrak L_s=\mathfrak K_r\mathfrak L_s\) es siempre de nuevo un anillo de matrices cuyo cuerpo de automorfismos \(\mathfrak M\) tiene por centro a \(\mathsf P\) ---el rango finito respecto de \(\mathsf P\) es evidente---. Por el teorema 1 de la página 20, \(\mathfrak K_r\times\mathfrak L_s\) es simple bilateral y completamente reducible, y es, por tanto, un anillo de matrices sobre un cuerpo de extensión \(\mathfrak M\) de \(\mathsf P\). Que el centro de \(\mathfrak M\) sea igual a \(\mathsf P\) se sigue inmediatamente de que, primero, \(\mathsf P\) está ciertamente contenido en el centro y, segundo, el centro ha de estar contenido en el centro \(\mathsf P\) de \(\mathfrak K\) (o de \(\mathfrak L\)). La multiplicación \(\mathfrak K_r\times\mathfrak L_s\) es manifiestamente asociativa. Los \(\mathfrak K_r\) forman un sistema cerrado multiplicativamente, pero no un grupo, pues si \(\mathfrak K_r\times\mathfrak L_s=\mathfrak M_m\), entonces ciertamente \(m\geq rs\); por consiguiente, no podrá hallarse ningún \(\mathfrak L_s\) que satisfaga la ecuación \(\mathfrak K_r\times\mathfrak L_s=\mathfrak M_m\) cuando \(r>m\).

Reunamos ahora en una clase \(\{\mathfrak K\}\) todos los \(\mathfrak K_r\) con el mismo \(\mathfrak K\). Entonces se cumple lo siguiente: si \(\mathfrak K_r\) y \(\mathfrak K_{r'}\) pertenecen a la misma clase \(\{\mathfrak K\}\), y \(\mathfrak L_s\) y \(\mathfrak L_{s'}\) a la misma clase \(\{\mathfrak L\}\), también \(\mathfrak K_r\times\mathfrak L_s\) y \(\mathfrak K_{r'}\times\mathfrak L_{s'}\) pertenecen a la misma clase. Basta mostrar que \(\mathfrak K_r\times\mathfrak L_s\) pertenece a la misma clase que \(\mathfrak K\times\mathfrak L\); es decir, que, si \(\mathfrak K\times\mathfrak L=\mathfrak M_m\), también \(\mathfrak K_r\times\mathfrak L_s\) es un anillo de matrices sobre \(\mathfrak M\). Pero esto se ve fácilmente: se tiene
\[
\mathfrak K_r\times\mathfrak L
=\left(\sum_1^r\mathfrak K c_{ik}\right)\mathfrak L
=\sum_1^r\mathfrak K_{\mathfrak L}c_{ik}
=\sum_1^r\mathfrak M_m c_{ik}.
\]
Este es un anillo de matrices de grado \(r\) sobre \(\mathfrak M_m\), y, por supuesto, un anillo de matrices de grado \(mr\) sobre \(\mathfrak M\). Prosiguiendo entonces exactamente del mismo modo,
\[
\mathfrak K_r\times\mathfrak L_s
=\left(\sum_1^s\mathfrak L d_{ik}\right)_{\mathfrak K_r}
=\sum_1^s\mathfrak L_{\mathfrak K_r}d_{ik}
=\sum_1^s\mathfrak M_{mr}d_{ik}
=\mathfrak M_{mrs},
\]
con lo cual queda demostrada nuestra afirmación. Retendremos que de \(\mathfrak L\times\mathfrak K=\mathfrak M_m\) se sigue \(\mathfrak K_r\times\mathfrak L_s=\mathfrak M_{mrs}\).

De lo demostrado se sigue que, si se define el producto de dos clases \(\{\mathfrak K\}\) y \(\{\mathfrak L\}\) como la clase del producto de un representante de \(\{\mathfrak K\}\) por un representante de \(\{\mathfrak L\}\), esta definición es independiente de la elección de los representantes.

Por consiguiente, la aplicación \(\mathfrak K_r\mapsto\{\mathfrak K\}\) es un homomorfismo del sistema de todos los \(\mathfrak K_r\) sobre el conjunto de todas las clases \(\{\mathfrak K\}\); el conjunto \(\mathscr K\) de las clases \(\{\mathfrak K\}\) es, pues, un sistema cerrado multiplicativamente. Demostraremos que \(\mathscr K\) es un grupo. Como se ve inmediatamente, su elemento unidad es la clase \(\{\mathsf P\}\), pues \(\mathfrak K_r\times\mathsf P=\mathfrak K_r\).

Solo queda, por tanto, demostrar para cada clase \(\{\mathfrak K\}\) la existencia de una clase \(\{\mathfrak K\}^{-1}\). Se encuentra que debemos poner \(\{\mathfrak K\}^{-1}=\{\overline{\mathfrak K}\}\), donde \(\{\overline{\mathfrak K}\}\) denota el cuerpo antiisomorfo a \(\mathfrak K\). Dicho de otro modo: \(\mathfrak K\times\overline{\mathfrak K}\) ha de ser un anillo de matrices sobre \(\mathsf P\); o bien, si \(\mathfrak K\) tiene rango \(t^2\), y por tanto \(\mathfrak K\times\overline{\mathfrak K}\) rango \(t^4\) sobre \(\mathsf P\), entonces \(\mathfrak K\times\overline{\mathfrak K}\) ha de descomponerse en \(t^2\) ideales izquierdos simples. Esto se sigue de que un ideal izquierdo simple de \(\mathfrak K\times\overline{\mathfrak K}\), considerado como módulo de representación de una representación recíproca irreducible de \(\overline{\mathfrak K}\) en \(\mathfrak K\) ---que es de grado 1---, tiene rango 1 sobre \(\mathfrak K\), y, por tanto, rango \(t^2\) sobre \(\mathsf P\).

Diremos que \(\mathfrak K\) es un \emph{cuerpo simple} si entre \(\mathsf P\) y \(\mathfrak K\) no existe ningún cuerpo cuyo centro sea también \(\mathsf P\).

\textbf{Teorema.} \emph{Todo cuerpo \(\mathfrak K\) es el producto directo de subcuerpos simples \(\mathfrak K^{(i)}\),}
\[
\mathfrak K=\mathfrak K^{(1)}\times\cdots\times\mathfrak K^{(p)}.
\]

\emph{Demostración.} Basta mostrar que, si \(\mathfrak K^{(1)}\) es un subcuerpo de \(\mathfrak K\) con centro \(\mathsf P\), existe un subcuerpo \(\mathfrak S\) de \(\mathfrak K\), también con centro \(\mathsf P\), que satisface la ecuación
\[
\mathfrak K^{(1)}\times\mathfrak S=\mathfrak K.
\]
En cualquier caso existe una clase \(\{\mathfrak S\}\) con la propiedad
\[
\{\mathfrak K^{(1)}\}\times\{\mathfrak S\}=\{\mathfrak K\},
\]
a saber,
\[
\{\mathfrak S\}=\{\overline{\mathfrak K^{(1)}}\}\times\{\mathfrak K\}.
\qquad
\text{Sea}\qquad
\overline{\mathfrak K^{(1)}}\times\mathfrak K=\mathfrak S_\lambda.
\]
Entonces
\[
\mathfrak K^{(1)}\times\mathfrak S_\lambda
=\mathfrak K^{(1)}\times\overline{\mathfrak K^{(1)}}\times\mathfrak K
=\mathsf P_{t^2}\times\mathfrak K
=\mathfrak K_{t^2},
\qquad
t^2=\text{rango de }\mathfrak K^{(1)}.
\]
Pero \(\mathfrak K^{(1)}\times\mathfrak S_\lambda\) es un anillo de matrices de grado al menos \(\lambda\), luego \(\lambda\leq t^2\). \(\mathfrak S_\lambda=\overline{\mathfrak K^{(1)}}\times\mathfrak K\) ha de descomponerse en ideales simples de rango 1 sobre \(\mathfrak K\), porque \(\overline{\mathfrak K^{(1)}}\) admite representaciones recíprocas de primer grado en \(\mathfrak K\); por consiguiente, \(\lambda\geq t^2\), \(\lambda=t^2\),
\[
\mathfrak K_{t^2}=\mathfrak K^{(1)}\times\mathfrak S_{t^2}.
\]
Por tanto, \(\mathfrak K^{(1)}\times\mathfrak S\) mismo ha de ser un cuerpo, y
\[
\mathfrak K^{(1)}\times\mathfrak S_{t^2}
=\bigl(\mathfrak K^{(1)}\times\mathfrak S\bigr)_{t^2}
=\mathfrak K_{t^2}.
\]
\(\mathfrak K^{(1)}\times\mathfrak S\) es el cuerpo de automorfismos de los ideales izquierdos simples de una descomposición de \(\mathfrak K_{t^2}\), y \(\mathfrak K\), el cuerpo de automorfismos de los ideales izquierdos simples de otra descomposición de \(\mathfrak K_{t^2}\). Pero, como cualesquiera dos ideales izquierdos distintos de \(\mathfrak K_{t^2}\) pueden incluirse en una misma descomposición, \(\mathfrak K\) y \(\mathfrak K^{(1)}\times\mathfrak S\) son isomorfos, y este isomorfismo puede elegirse de modo que deje fijo elemento a elemento a \(\mathsf P\). Por ello, según el teorema 7 de la página 25, este isomorfismo está engendrado por un automorfismo interior de \(\mathfrak K_{t^2}\).
\[
\mathfrak K
=\lambda^{-1}\bigl(\mathfrak K^{(1)}\times\mathfrak S\bigr)\lambda
=\lambda^{-1}\mathfrak K^{(1)}\lambda\times\lambda^{-1}\mathfrak S\lambda,
\qquad
\lambda\in\mathfrak K_{t^2}.
\]
\(\lambda^{-1}\mathfrak K^{(1)}\lambda\) es un subcuerpo de \(\mathfrak K\) isomorfo a \(\mathfrak K^{(1)}\); por tanto, según el teorema 7 de la página 25, para un \(\mu\) conveniente de \(\mathfrak K\),
\[
\mu^{-1}\lambda^{-1}\mathfrak K^{(1)}\lambda\mu=\mathfrak K^{(1)},
\]
y, por consiguiente,
\[
\mu^{-1}\mathfrak K\mu
=\mathfrak K
=\mathfrak K^{(1)}\times\mu^{-1}\lambda^{-1}\mathfrak S\lambda\mu,
\]
con lo cual queda demostrada la afirmación.

\subsection*{§ 23. Sistemas de factores}\label{faktorensysteme}

En las investigaciones que siguen se supondrá que el cuerpo base \(\mathsf P\) es perfecto. Esta hipótesis puede sustituirse por la más general de considerar únicamente aquellos cuerpos no conmutativos \(\mathfrak K\) con centro \(\mathsf P\) que poseen la propiedad de que todo \(\mathfrak K_r\) contiene un subcuerpo conmutativo maximal que es de primera especie sobre \(\mathsf P\). (Como Köthe ha demostrado entretanto, esto se cumple para todo \(\mathfrak K\).)

Sea \(Z\) un cuerpo de descomposición del cuerpo no conmutativo \(\mathfrak K\) con centro \(\mathsf P\), y sea \(\mathfrak Z\) la representación irreducible de \(Z\) en \(\mathfrak K\), de grado \(r\); es entonces un subcuerpo conmutativo maximal de \(\mathfrak K_r\). Sea \(\Gamma\) el cuerpo de Galois de \(Z\); por tanto,
\[
\mathfrak Z_\Gamma=e_1\Gamma+\cdots+e_n\Gamma
\qquad\text{(página 9)}.
\]
Aquí \(n\) es el grado de \(Z\) sobre \(\mathsf P\). Para el anillo \(\mathfrak K_r\) extendido mediante \(\Gamma\) se obtiene asimismo una descomposición
\[
\mathfrak K_{r\Gamma}=\mathfrak K_{r\Gamma}e_1+\cdots+\mathfrak K_{r\Gamma}e_n,
\]
\[
\mathfrak K_{r\Gamma}=\mathfrak l_1+\cdots+\mathfrak l_n,
\qquad
\mathfrak l_i=\mathfrak K_{r\Gamma}e_i.
\]
Los \(\mathfrak l_i\) son manifiestamente ideales izquierdos. Son incluso ideales izquierdos simples, pues, como \(Z\) es un cuerpo de descomposición, \(\mathfrak K_{r\Gamma}\) es un anillo de matrices de grado \(r\cdot t=n\) sobre \(\Gamma\). (Teorema 4, página 23; \(n=rt\) en la página 22.) La descomposición así obtenida de \(\mathfrak K_{r\Gamma}\) en ideales izquierdos simples está especialmente adaptada a las propiedades algebraicas de \(\mathfrak K\) y \(Z\). En efecto, si, como en la página 18, se define la acción de una sustitución \(S\) del grupo de Galois \(\mathfrak G\) de \(\Gamma/\mathsf P\) sobre todos los elementos de \(\mathfrak K_{r\Gamma}\) poniendo, para cada elemento \(z\) de \(\mathfrak K_r\),
\[
S(z)=z,
\]
entonces los ideales \(\mathfrak l_i\) son conjugados; esto es, la aplicación de un \(S\) a un \(\mathfrak l_j\) da otro,
\[
S(\mathfrak l_j)=\mathfrak l_i,
\]
porque toda componente indescomponible \(e_j\) de la unidad debe pasar a otra \(e_i\). (Véase también la página 11.)

El módulo de representación \(e_i\Gamma\) aplica \(\mathfrak Z\) sobre un subcuerpo \(Z_i\) de \(\Gamma\). El elemento \(e_i\) pertenece ya a \(\mathfrak Z_{Z_i}\), porque la representación \(Z_i\) de \(\mathfrak Z\) está ya contenida en \(Z_i\), de modo que \(\mathfrak Z_{Z_i}\) ha de separar el módulo de representación \(e_iZ_i\). Sea \(\{Z_i,Z_k\}\) el cuerpo compuesto de \(Z_i\) y \(Z_k\). \(\mathfrak K_{r\{Z_i,Z_k\}}\) separa los dos ideales izquierdos
\[
\bar{\mathfrak l}_i=\mathfrak K_{r\{Z_i,Z_k\}}e_i,
\qquad
\bar{\mathfrak l}_k=\mathfrak K_{r\{Z_i,Z_k\}}e_k,
\]
que son simples, porque \(\mathfrak l_i=\mathfrak K_{r\Gamma}\bar{\mathfrak l}_i\). Por tanto,
\[
\bar{\mathfrak l}_i=\sum_{\lambda=1}^m\{Z_i,Z_k\}c_{\lambda i}^{(ik)},
\qquad
\bar{\mathfrak l}_k=\sum_{\lambda}\{Z_i,Z_k\}c_{\lambda k}^{(ik)}.
\]
Aquí \(c_{\lambda\rho}^{(ik)}\) es un sistema de unidades matriciales de \(\mathfrak K_{r\{Z_i,Z_k\}}\).

Todo isomorfismo \(a_i\mapsto a_k\) de \(\bar{\mathfrak l}_i\) sobre \(\bar{\mathfrak l}_k\), con \(\mathfrak K_{r\{Z_i,Z_k\}}\) como dominio de operadores, queda caracterizado por el elemento \(p_{ik}\) mediante \(e_i\mapsto p_{ik}\), pues entonces \(a_i=a_ie_i\mapsto a_ip_{ik}\), y, por tanto, \(\bar{\mathfrak l}_i\cdot p_{ik}=\bar{\mathfrak l}_k\). Todos los \(p_{ik}\) posibles de esta clase son, manifiestamente,
\[
p_{ik}=c_{ik}^{(ik)}\gamma_{ik},
\qquad
\gamma_{ik}\ne0\quad\text{de}\quad \{Z_i,Z_k\}.
\]
La aplicación \(a_i\mapsto a_ip_{ik}\) proporciona entonces también un isomorfismo de operadores de \(\mathfrak l_i\) sobre \(\mathfrak l_k\). Elegiremos ahora sistemas determinados de tales \(p_{ik}\), imponiéndoles una condición de conjugación.

Definamos la acción de un \(S\) sobre un índice \(i\) mediante
\[
S(i)=k,
\qquad\text{si}\qquad
S(Z_i)=Z_k.
\]
Para cada par de índices \((\nu,\mu)\) existe al menos un par conjugado \(S(\nu,\mu)\) de la forma \(S(\nu,\mu)=(1,k)\). De cada clase de pares de índices conjugados se elige un par fijo \((1,k)\), y para este \((1,k)\) se pone
\[
p_{1k}=c_{1k}^{(1k)}\gamma_{1k},
\qquad
\gamma_{1k}\ne0\quad\text{de}\quad \{Z_1,Z_k\},
\]
de manera por lo demás arbitraria; si \(S(1,k)=(\nu,\mu)\), se pone entonces
\[
p_{\nu\mu}=S(p_{1k}).
\]
En efecto, \(p_{\nu\mu}\) tiene entonces la forma
\[
p_{\nu\mu}=c_{\nu\mu}^{(\nu\mu)}\gamma_{\nu\mu},
\qquad
\gamma_{\nu\mu}\ne0\quad\text{de}\quad \{Z_\nu,Z_\mu\}.
\]
Pues la acción de \(S\) sobre el isomorfismo
\[
a_1\mapsto a_1p_{1k},
\qquad\text{en particular}\qquad
e_1\mapsto p_{1k}
\quad\text{de }\bar{\mathfrak l}_1\text{ sobre }\bar{\mathfrak l}_k
\]
da un isomorfismo
\[
a_\nu\mapsto a_\nu p_{\nu\mu},
\qquad\text{en particular}\qquad
e_\nu\mapsto p_{\nu\mu}
\quad\text{de }\bar{\mathfrak l}_\nu\text{ sobre }\bar{\mathfrak l}_\mu.
\]

Si por \(c_{\lambda\rho}\) se entiende un sistema de unidades matriciales de \(\mathfrak K_{r\Gamma}\), es claro que \(p_{\nu\mu}\) puede escribirse también en la forma
\[
p_{\nu\mu}=c_{\nu\mu}\bar\gamma_{\nu\mu}.
\]
Las relaciones
\[
c_{ij}c_{1k}=0,\quad\text{si}\quad j\ne1,
\qquad
c_{ij}c_{jk}=c_{ik}
\]
tienen, por consiguiente, como consecuencia las relaciones
\[
\begin{aligned}
p_{ij}p_{1k}&=0,
&&\text{si }j\ne1,\\
p_{ij}p_{jk}&=\alpha_{ik}^{(j)}p_{ik},
&&\alpha_{ik}^{(j)}\text{ de }\Gamma
\text{ (más precisamente, de }\{Z_i,Z_j,Z_k\}\text{)}.
\end{aligned}
\]
Las \(n^3\) magnitudes \(\alpha\) forman el \emph{sistema de factores de \(\mathfrak K\) respecto de \(Z\)}, asociado al sistema \(p_{ik}\) de ``pseudounidades matriciales''.

Las \(\alpha\) son conjugadas: de \(S(i,k,j)=(\nu,\mu,\tau)\) se sigue
\[
S\bigl(\alpha_{ik}^{(j)}\bigr)=\alpha_{\nu\mu}^{\tau}.
\]

Todo sistema \(p_{ik}^*\) de pseudounidades matriciales se obtiene de un sistema \(p_{ik}\) multiplicando los \(p_{ik}\) iniciales por elementos arbitrarios \(\delta_{ik}\ne0\) de \(\{Z_i,Z_k\}\):
\[
p_{1k}^*=p_{1k}\delta_{1k},
\]
y poniendo de nuevo, cuando \(S(1,k)=(\nu,\mu)\),
\[
p_{\nu\mu}^*=S(p_{1k}^*).
\]
Por tanto,
\[
p_{\nu\mu}^*=p_{\nu\mu}\delta_{\nu\mu},
\]
con \(\delta_{\nu\mu}\) de \(\{Z_\nu,Z_\mu\}\); los \(\delta_{\nu\mu}\) son conjugados, y de \(S(i,k)=(\nu,\mu)\) se sigue \(S(\delta_{ik})=\delta_{\nu\mu}\). Si \(\alpha^*\) es el sistema de factores asociado a \(p_{ik}^*\), se tiene
\[
\alpha_{ik}^{(j)*}
=\frac{\delta_{ij}\delta_{jk}}{\delta_{ik}}\alpha_{ik}^{(j)}.
\]

Los sistemas de factores \(\alpha^*\) y \(\alpha\) se llaman asociados (y también \(p_{ik}^*\) y \(p_{ik}\)); todos los sistemas de factores de \(\mathfrak K\) respecto de \(Z\) forman una clase \(\{\alpha\}\) de sistemas de factores asociados.

La clase \(\{\alpha\}\) no depende de la representación \(\mathfrak Z\) de \(Z\) en \(\mathfrak K_r\) que se tome como base de su definición. En efecto, como el paso de un \(\mathfrak Z\) a otro puede realizarse mediante transformación por un elemento \(\lambda\) de \(\mathfrak K_r\), que es, por tanto, \(\mathfrak G\)-invariante, de un sistema \(p_{ik}\) asociado a \(\mathfrak Z\) se obtiene un sistema
\[
p'_{ik}=\lambda^{-1}p_{ik}\lambda,
\]
asociado a \(\mathfrak Z'\); pero los sistemas de factores de \(p_{ik}\) y \(p'_{ik}\) son idénticos.

Un sistema de factores no es sino una tabla de multiplicación de \(\mathfrak K_{n\Gamma}\), pero una tabla adaptada a las propiedades algebraicas particulares de \(\mathfrak K\).
